\subsection{General original Gauss comparison helpers}

\begin{lemma}[Upper bound from the original central fibre only]
\label{lem:nf-general-gauss-usc-upper}
\lean{EnvelopeProperty.eventually_flatNormFun_nonflatGaussPoint_lt}
\leanok
\uses{lem:nf-pieces,lem:nf-gauss-ray}
Let $k$ be any field and let $F$ be an algebraically closed ultrametric
normed field with its original $k$-algebra structure. Assume
$\|a\|_F=1$ for every nonzero $a\in k$ and that $F$ has an element of
norm greater than one. Let $p\in F[T]$ satisfy $p(0)\ne0$, put
$Q=F[T,1/p]$, and retain the original evaluation point $\xi$ and
Gauss points $\zeta_r$ in $\operatorname{Berk}_k Q$.
Let $R$ be ANY commutative $Q$-algebra, finite free over $Q$, with
its original compatible $k$-action. For any original function
$U:\operatorname{Berk}_k R\to\overline{\mathbb R}$, assume only that
$U$ is upper semicontinuous at EVERY point over $\xi$ and that
$U(x)\ne+\infty$ at each such point. Then for EVERY
$C\in\overline{\mathbb R}$ with $N_{R/Q}U(\xi)<C$,
\[
 N_{R/Q}U(\zeta_r)<C
 \quad\text{for all sufficiently small }r\in\mathbb R_{\ge0}.
\]
Here the norm is the literal original finite sum with its completed
local-fibre multiplicities. No FS presentation, local psh class,
regularizability, completeness of $F$, domain or product decomposition
of $R$, or nonempty-fibre premise is assumed.
\end{lemma}
\begin{proof}
Use the existing genuine finite-fibre pieces near $\xi$. Their actual
positive weights have the same total mass over each central point for
all sufficiently small Gauss radii. Since the central values exclude
$+\infty$ and their weighted sum is strictly less than $C$, choose
real cuts $c_x>U(x)$ whose finite weighted sum is still strictly less
than $C$. This also applies when some central values are $-\infty$;
when the central fibre is empty the norm is the empty sum.
Upper semicontinuity gives an original neighbourhood of each central
point on which $U<c_x$. The existing shrinking and finite synchronization
put each corresponding Gauss piece in that neighbourhood. Sum these
inequalities with the SAME nonnegative original weights and use the
proved mass equality. The resulting sum is at most the weighted sum
of the chosen real cuts, hence strictly below $C$. No lower-bound or
approximation theorem is used.
\end{proof}

\begin{lemma}[Whole order squeeze along positive radii]
\label{lem:nf-general-order-squeeze}
\lean{EnvelopeProperty.nonflat_eq_of_squeeze}
\leanok
Let $N,V:\mathbb R_{\ge0}\to\overline{\mathbb R}$ and
$N_0,v_0\in\overline{\mathbb R}$. Assume that eventually as $r\to0$,
$N(r)=V(r)$ whenever $r>0$, that eventually $N_0\le N(r)$ and
$v_0\le V(r)$, and that for EVERY extended-real $C>N_0$
respectively $C>v_0$ one eventually has $N(r)<C$ respectively
$V(r)<C$. Then $v_0=N_0$. There is no geometric, function-class,
finite-valuedness, countability, or field hypothesis.
\end{lemma}
\begin{proof}
For every $C>N_0$, intersect the eventual upper bound for $N$, the
lower bound for $V$, and equality on positive radii, in the punctured
right neighbourhood filter of zero. This filter is nontrivial, so
choose an actual positive radius in the intersection. It gives
$v_0\le N(r)<C$, and the order-density criterion implies $v_0\le N_0$.
Interchanging the roles gives $N_0\le v_0$. Antisymmetry proves the
whole equality, including either infinite endpoint.
\end{proof}

\begin{lemma}[The entire rigid norm identity for every original function]
\label{lem:nf-whole-rigid-weighted-norm}
\lean{EnvelopeProperty.NonflatRigidEmbedding.lineWeights,EnvelopeProperty.NonflatRigidEmbedding.flatNormFun_line_eq_sum_rigidWeight}
\leanok
\uses{lem:nf-line-weights,def:nf-weights,algebra:finite-free-normalized-line-product,algebra:whole-normalized-line-product-map}
Let $k$ be algebraically closed, let $A,B$ be original commutative
domains over $k$, and let $\varphi:B\to A$ be finite and injective.
Retain a nonzero principal-open presentation $B\simeq k[x_1,\ldots,x_n]_g$,
an actual rigid datum $\tau:H(s)\hookrightarrow F$ for an original
$s\in\operatorname{Berk}_k B$, with $F$ algebraically closed
ultrametric normed and its original norm trivial on $k$, and the actual
full good locus $h$ of $\varphi$ (including its original FREE and ETALE
fields). Let $w$ be an actual direction, with a specified nonzero
coordinate, admissible for $h$. Use the SAME resulting open line $Q$,
normalized product $R'$, and maps $r,p$.
For an original separating element $b\in F\otimes_k A$, let $M_z$
be its actual characteristic-polynomial root multiplicities at the
rigid fibre, and let $m_y=\sum_{r(z)=y}M_z$ be the existing rigid weights.
The EXISTING whole line theorem gives
$\sum_z M_z=\operatorname{genericDegree}(\varphi)$ and, for EVERY
$W:\operatorname{Berk}_k(F\otimes_k A)\to\overline{\mathbb R}$,
\[
 N_{R'/Q}(W\circ p)(\xi)=\sum_zM_zW(z).
\]
Consequently for EVERY original $u:\operatorname{Berk}_k A\to\overline{\mathbb R}$,
\[
 N_{R'/Q}(u\circ r\circ p)(\xi)
   =\sum_{\varphi^*y=s}m_yu(y).
\]
There is no psh, regularizability, USC, growth, countability, Wu or
completeness hypothesis. Every coefficient, point map and local-fibre
multiplicity is the original one.
\end{lemma}
\begin{proof}
The existing whole line-weight theorem uses the genuine original finite-etale good locus and the separating element to give the displayed degree mass and the entire equality for every $W$. Specialize it to $W=u\circ r$ to obtain $\sum_zM_zu(r(z))$. The actual rigid fibre is finite,
and every original $z$ restricts to a point $y$ of the original fibre
over $s$. Regroup this actual finite sum by that restriction map.
The existing definition of $m_y$ and finite distributivity for the
nonnegative integer weights give exactly the displayed right side.
This is the original whole-function identity for arbitrary extended-real
values; no approximation or change of weights enters it.
\end{proof}
